
\setcounter{equation}{0}
\setcounter{Q}{0}

Le diagramme binaire isobare solide-liquide bismuth-étain est représenté sur la figure \ref{}, sous une pression $P = P^\circ = 1 ~ \si{\bar}$, avec en abscisse la fraction massique en étain, $w(\ce{Sn})$ et en ordonnée la température exprimée en degré Celsius. L'étain qui possède une structure
cubique en dessous de \SI{13}{\celsius}, puis quadratique au-dessus, et le bismuth (de structure hexagonale) ne sont que partiellement miscibles à l'état solide. $\alpha$ et $\beta$ désignent des solutions solides, respectivement de \ce{Sn} dans \ce{Bi} et de \ce{Bi} dans \ce{Sn}. Toutes les notations et les calculs
feront exclusivement référence aux fractions massiques des constituants.

\begin{center}
\begin{minipage}{0.58\textwidth}
\includegraphics[width=\textwidth]{/home/jeff/Documents/IUTOrsay/1-Cours/4p13-Alliages/Cours-TD/Exos/Binaires-Lecture/MiscibilitePartielle-Ex01/figures/Fig01}
\end{minipage}\begin{minipage}{0.38\textwidth}
\captionof{figure}{Diagramme binaire isobare, à une pression de \SI{1}{\bar}, du mélange \ce{Bi}-\ce{Sn}.}
[La composition est donnée par la fraction massique en étain dans le mélange]
\end{minipage}
\end{center}

\Q Indiquer le nombre et la nature des phases en présence dans les différentes régions du diagramme binaire.

\Q Rappeler les définitions du solidus et du liquidus ; préciser leur localisation sur le diagramme.

\Q Expliquer pourquoi la miscibilité de \ce{Sn} et \ce{Bi} ne peut pas être totale à l'état solide.

\Q À quoi correspondent les points notés \textbf{A} et \textbf{B} sur le diagramme ? Que représentent les abscisses des points notés \textbf{C} et \textbf{D} ?

\Q Nommer le point \textbf{E}. Écrire l'équilibre existant en ce point et déterminer sa variance particulière.


Préparons un mélange contenant \SI{10}{\kilo\gram} d'étain et \SI{90}{\kilo\gram} de bismuth ; le mélange est fondu, homogénéisé puis refroidi lentement.

\Q Déterminer :
\begin{itemize}
\item la fraction massique $w(\ce{Sn})$ de cet alliage ;
\item la température de début de solidification du mélange ;
\item la teneur approximative en étain du premier cristal qui se forme ;
\item les masses et la composition des phases à \SI{150}{\celsius}.
\end{itemize}

\paragraph*{Étude de l'alliage 42Sn-58Bi (\SI{42}{\percent} massique d'étain et \SI{58}{\percent} massique de bismuth)}

\Q Préciser l'intérêt d'utiliser un alliage de cette composition.

Le mélange fondu est refroidi depuis \SI{250}{\celsius} jusqu'à \SI{100}{\celsius}, où il subit une trempe (refroidissement brutal) jusqu'à la température ambiante.

\Q Tracer l'allure de la courbe de refroidissement de cet alliage en mentionnant les températures de rupture de pente et la nature des phases en présence durant chaque étape du processus de solidification jusqu'à \SI{100}{\celsius}.

\Q Recenser les phases en présence et indiquer leur composition à \SI{100}{\celsius}.

Lors du refroidissement, des prélèvements de matière sont effectués puis subissent une trempe ; après polissage, ils peuvent être observés à l'aide d'un microscope métallographique afin de vérifier leur microstructure. Les clichés microstructuraux d'un échantillon de ce mélange ayant subi un refroidissement lent de \SI{250}{\celsius} à la température ambiante sont schématisés (dans le désordre) sur la figure ci-dessous, où les symboles $\alpha$, $\beta$ et $L$ sont relatifs aux diverses phases en présence.

\begin{figure}[H]
\begin{minipage}{0.5\textwidth}
\centering
\subfloat[\label{fig:MP-Ex01-02a}]{\includegraphics[width=2cm]{/home/jeff/Documents/IUTOrsay/1-Cours/4p13-Alliages/Cours-TD/Exos/Binaires-Lecture/MiscibilitePartielle-Ex01/figures/Fig02b}}
\hspace{0.5cm}\subfloat[\label{fig:MP-Ex01-02b}]{\includegraphics[width=2cm]{/home/jeff/Documents/IUTOrsay/1-Cours/4p13-Alliages/Cours-TD/Exos/Binaires-Lecture/MiscibilitePartielle-Ex01/figures/Fig02a}}
\hspace{0.5cm}\subfloat[\label{fig:MP-Ex01-02c}]{\includegraphics[width=2cm]{/home/jeff/Documents/IUTOrsay/1-Cours/4p13-Alliages/Cours-TD/Exos/Binaires-Lecture/MiscibilitePartielle-Ex01/figures/Fig02c}}
\end{minipage}\hfill\begin{minipage}{0.5\textwidth}
\caption{Schéma des clichés de miscrocopie optique obtenu en différents points (compositions).}
\label{fig:MP-Ex01-02}
\end{minipage}
\end{figure}

\Q Expliquer, en le justifiant, à quelle température (ou dans quel intervalle de température) chacun de ces clichés, notés respectivement \ref{fig:MP-Ex01-02a}, \ref{fig:MP-Ex01-02b} et \ref{fig:MP-Ex01-02c}, peut être observé au microscope.